\chapter{Differences from C GLoBES}\label{ch:cdiffs}

GLoBES.jl was developed as a port of C GLoBES~3.0.11 and is tested against it.
Where the C library has a bug, GLoBES.jl does the right thing instead. This
chapter lists every such difference, how it shows, and how much it changes
results. The test suite compares GLoBES.jl with a C GLoBES built with the
patch \file{test/reference/c-fixes.patch}, which applies the same fixes to the
C code where that is a small change (Sections~\ref{sec:cdiff-constants},
\ref{sec:cdiff-flux} and~\ref{sec:cdiff-tilt}). With it, both agree to
rounding errors except where noted.

\section{Physical constants}\label{sec:cdiff-constants}

\begin{center}\small
\begin{tabular}{lll}
\toprule
& C GLoBES & GLoBES.jl (CODATA 2018) \\
\midrule
$\hbar c$ & $1.974\,7235\times10^{-10}\eV\km$ & $1.973\,269\,804\times10^{-10}\eV\km$ \\
$\sqrt2 G_F N_A$ & $7.5\times10^{-14}\eV$ per $\mathrm{g/cm^3}$ & $7.6325\times10^{-14}\eV$ per $\mathrm{g/cm^3}$ \\
$\pi$ (in fluxes) & $3.1415$ & $\pi$ \\
$N_A$ (flux normalisation) & $6.02204\times10^{23}$ & $6.022\,140\,76\times10^{23}$ \\
$m_\mu$, $m_e$ & $0.1056$, $0.511\times10^{-3}\GeV$ & $0.105\,658\,3755$, $0.510\,998\,95\times10^{-3}\GeV$ \\
\bottomrule
\end{tabular}
\end{center}
The first shifts every oscillation phase by $0.07\%$ (equivalent to a change of
$\dm{}$ by that amount), the second makes the matter potential $1.7\%$ larger,
which is well within the usual density uncertainty but noticeable in precise
comparisons. The others change built-in fluxes by $10^{-5}$ to $10^{-4}$.

\section{Flux tables}\label{sec:cdiff-flux}

A flux table has 501 equidistant energies $E_0,\dots,E_{500}$. To interpolate,
C GLoBES computes the step as $(E_{500}-E_0)/501$ instead of $/500$, so it
reads the table at the wrong position: every flux is evaluated at $1.002$
times the requested energy (more precisely at $E_0+\frac{501}{500}(E-E_0)$),
and the last $0.2\%$ of the table is never used. For the built-in fluxes, which
are tabulated at initialisation, it also places the 501 points up to $500/501$
of the parent energy, so that the highest-energy part of the neutrino factory
spectrum, where the $\numu$ flux is largest, is cut off. GLoBES.jl uses the
correct step and tabulates the built-in fluxes up to the kinematic endpoint.
The beta beam flux uses the phase space integral $g(y_e)$ with $\log_{10}$
instead of $\ln$, which changes its normalisation by about $0.35\%$.

Together with the constants above, this changes total event numbers by
$0.1$--$1\%$ for the superbeam, reactor and beta beam experiments of
\file{test/data}, and by up to $3.5\%$ for the neutrino factory appearance
channels near the end of the spectrum.

\section{Background tilt prior}\label{sec:cdiff-tilt}

In \code{chiSpectrumTilt} and \code{chiTotalRatesTilt}, the nuisance parameter
of the background tilt $\zeta_4$ enters the rates as $\zeta_4/\Delta E$, like
the signal tilt. C GLoBES applies the prior to $\zeta_4/\Delta E$, while the
prior of the signal tilt is on $\zeta_2$:
\[
\text{C:}\quad \frac{\zeta_2^2}{\sigma_2^2} + \frac{(\zeta_4/\Delta E-\tau_0)^2}{\sigma_4^2},
\qquad
\text{GLoBES.jl:}\quad \frac{\zeta_2^2}{\sigma_2^2} + \frac{(\zeta_4-\tau_0)^2}{\sigma_4^2}.
\]
In C the background tilt error is therefore $\sigma_4\,\Delta E$ instead of the
$\sigma_4$ written in the AEDL file, a factor of $0.8$ for T2K and $46$ for the
neutrino factory. (In \code{chiSpectrumCalib} both priors are on the
parameters.)

\section{Empty bins in the Poisson \texorpdfstring{$\chi^2$}{chi2}}\label{sec:cdiff-poisson}

C GLoBES returns $10^{100}$ from \code{glb\_likelihood} whenever the prediction
is $N\le0$, even for a bin with no data, $n=0$. Experiments with empty bins
(for example a beta beam or neutrino factory with bins above the beam energy,
\file{NF\_hR\_lT.glb} and \file{BBvar\_WC.glb} in \file{test/data}) then have
$\chi^2=10^{100}$ for every hypothesis, which is meaningless; and when
systematics push a prediction through zero, $\chi^2$ jumps from a finite value
to $10^{100}$, which disturbs the minimisers. GLoBES.jl gives empty bins
$\chi^2_\text{P}=2|N|$ and continues the Poisson $\chi^2$ smoothly below
$N=10^{-10}$ (Section~\ref{sec:poisson}).

\section{Energy scale shift}\label{sec:cdiff-escale}

C GLoBES applies an energy calibration error $b$ (in \code{chiSpectrumCalib}
and \code{glbShiftEnergyScale}) by linear interpolation of the bin contents at
the shifted positions $\delta_i=b(i+t_0+\frac12)+i$:
$\mathcal S_b[s]_i=(1+b)\,[s_k+(s_{k+1}-s_k)(\delta_i-k)]$, $k=\lfloor\delta_i\rfloor$.
This has three problems:
\begin{enumerate}
\item The interpolation is between bin \emph{contents}. Since $\delta_i$ moves
  across bin boundaries at $b=0$, the result has a kink at $b=0$ in every bin:
  $\chi^2(b)$ has a cusp at its minimum (Figure~\ref{fig:escale}), and its
  curvature, which determines the effect of the error, is ill-defined.
\item For $k=N_\text{bins}-1$ (shifts beyond the upper end), it uses
  $s_k(\delta_i-k)+s_k$, which counts the last bin twice instead of assuming
  zero events above the range.
\item It assumes bins of equal width; with \code{\$binsize} the shift is wrong.
\end{enumerate}
GLoBES.jl integrates the monotone interpolation of the cumulative spectrum
between the scaled bin edges (Section~\ref{sec:escale}), which avoids all three.
The test suite therefore uses \code{chiSpectrumTilt} instead of
\code{chiSpectrumCalib} in its comparisons with C, and tests the new shift
against exact integrals of known spectra.

\section{The minimiser}\label{sec:cdiff-powell}

Powell's method replaces the direction of largest decrease by the new
direction if~\cite{NR}
\begin{equation}
t = 2(f_0-2f_N+f_E)\,(f_0-f_N-\Delta f)^2 - \Delta f\,(f_0-f_E)^2 < 0 .
\end{equation}
C GLoBES writes this with the Numerical Recipes macro
\code{SQR(a)}, defined as \code{(sqrarg=(a), sqrarg*sqrarg)} with a single
static variable \code{sqrarg}. Two \code{SQR} in one expression are undefined
behaviour; GCC evaluates both assignments before both products and computes
$t=(2(f_0-2f_N+f_E)-\Delta f)(f_0-f_E)^2$. The test then accepts or rejects
new directions for the wrong reasons. In addition, a single Powell run can stop
early in a long curved valley, as around degenerate solutions. GLoBES.jl uses
the formula above and restarts the minimisation from the result with fresh
directions until a whole run no longer improves $\chi^2$ (at most ten times).
Fits may therefore end at a lower $\chi^2$ than in C GLoBES.

\section{Data file lookup}

C GLoBES looks for flux and cross section files only in the current directory
and in \code{GLB\_PATH}. GLoBES.jl also looks next to the AEDL file, so that
\code{Experiment("/path/to/exp.glb")} works from any directory.

\section{State and interface}

\begin{itemize}
\item All state belongs to a \code{Setup} instead of global variables
  (Section~\ref{sec:setups}).
\item The high-level API raises exceptions instead of printing errors and
  returning codes; the \code{glb*} functions keep the C behaviour.
\item Numbering starts at 1 (Chapter~\ref{ch:compat}).
\item The rates of a rule outside its energy window are zero (C leaves them
  uninitialised).
\end{itemize}
